% !TeX root = main.tex
\section{工业机器人I}

\subsection{机器人结构与运动学}

\paragraph{机器人的分类}

在工厂里，我们经常需要设备把零件抓起来、移动、放下，这类设备统称为操作设备（Handhabungsgeräte）。它们大致分三种：
\begin{itemize}
    \item 靠人手动操作或远程遥控的装置
    \item 固定程序运动的装置
    \item 自由可编程运动的装置
\end{itemize}
工业机器人是一种自动控制的、可重复编程的多用途操作机，可在三个或三个以上轴上编程，既可固定安装，也可移动使用，用于工业自动化应用。属于自由可编程运动的装置。

工业机器人本身还可以按其程序自主性的高低做进一步细分，这关注的是机器人在运行过程中能在多大程度上自己做主：
\begin{itemize}
    \item \textbf{无自主程序影响（ohne selbsttätige Programmbeeinflussung）}：机器人按照固定程序执行。
    \item \textbf{有自主程序选择（mit selbsttätiger Programmselektion）}：机器人事先存有多套可选程序，运行时能够自己判断该调用哪一套程序。
    \item \textbf{有自主程序适应（mit selbsttätiger Programmadaption）}：机器人能根据实时感知到的环境或工件状态，对程序本身的参数甚至逻辑做出动态调整。
\end{itemize}

\paragraph{工业机器人的结构}

工业机器人可以看作由多个刚性连杆和关节组成的运动链。连杆（Glied）是机器人中近似刚性的构件。关节（Gelenk）连接两个连杆并允许它们产生相对运动的。轴（Achse）是指机器人中被驱动和控制的运动轴。关节的运动主要分为转动和平移。多个关节和连杆依次连接，构成了机器人的运动链。

自由度（Freiheitsgrad）表示一个机械系统能够独立运动的参数数量。关节数量不一定等于自由度数量。因为不同关节产生的运动可能存在约束，并非每一个关节都提供一个独立自由度。

\paragraph{坐标变换}

机器人系统中会同时存在多个坐标系，同一个点在不同坐标系中的坐标不同，因此需要通过\textbf{坐标变换}建立他们之间的关系。

假设有两个坐标系 \(A\) 和 \(B\)。如果 \(B\) 相对于 \(A\) 只是移动了，没有旋转，那么一个点从 \(B\) 坐标系转换到 \(A\) 坐标系时，只需要加上两个坐标系原点之间的位移：
\begin{equation}
    {}^A\mathbf p={}^B\mathbf p+{}^A\mathbf t_B
\end{equation}

但现实中 \(B\) 往往还相对于 \(A\) 发生了旋转。因此首先需要用旋转矩阵 \({}^A R_B\) 把向量的方向从 \(B\) 的坐标轴转换成 \(A\) 的坐标轴：

\begin{equation}
    \label{eq:rotationtranslation}
    {}^A\mathbf p={}^A R_B\,{}^B\mathbf p+{}^A\mathbf t_B
\end{equation}

这里 ${}^A R_B$ 表示“坐标系 \(B\) 相对于坐标系 \(A\) 的方向关系”，而 ${}^A\mathbf t_B$ 表示“\(B\) 的原点在 \(A\) 中的位置”。可以记成：下标是从哪里来，上标是到哪里去。

但是表达式 \ref{eq:rotationtranslation} 还是有一点麻烦，因为旋转是乘矩阵，平移却是加向量。为了让旋转和平移能够统一用一次矩阵乘法完成，我们引入第四个齐次坐标，把三维点 \((x,y,z)^T\) 写成 \((x,y,z,1)^T\)。然后定义齐次变换矩阵为
\begin{equation}
    {}^A T_B=
    \begin{bmatrix}
        {}^A R_B & {}^A\mathbf t_B\\
        0&0&0&1
    \end{bmatrix}
\end{equation}

这样表达式 \ref{eq:rotationtranslation} 就可以写成
\begin{equation}
    \begin{bmatrix}
        {}^A\mathbf p\\
        1
    \end{bmatrix}
    ={}^A T_B
    \begin{bmatrix}
        {}^B\mathbf p\\
        1
    \end{bmatrix}
\end{equation}

这就是齐次变换的意义：用一个矩阵同时表示两个坐标系之间的旋转和平移。

\paragraph{示例：坐标变换}

图 \ref{fig:roboterundwelt} 展示了一个机器人在操控一个物体。这个场景有4个坐标系，分别是世界坐标系W、机器人的基座坐标系B、机器人末端坐标系TCP、以及物体坐标系O。图中已知三个变换 \({}^W T_B\)、\({}^B T_{TCP}\)、和\({}^W T_O\)；两个未知变换 \({}^W T_{TCP}\) 和 \({}^{TCP}T_O\)。首先我们要明确：\({}^W T_B\) 表示 \(B\to W\)。接下来我们要来求这两个未知变换。

\begin{figure}[htbp]
  \centering
  \includegraphics[width=0.5\textwidth]{roboterundwelt.png}
  \caption{一个机器人在环境中操控一个物体，它们有各自的坐标。}
  \label{fig:roboterundwelt}
\end{figure}

我们使用坐标变换的串联将 \(TCP\to B\) 和 \(B\to W\) 连接起来，自然就得到了 \(TCP\to B\to W\)。对应的矩阵就是 \({}^W T_B\,{}^B T_{TCP}\)。因此
\begin{equation}
{}^W T_{TCP}={}^W T_B\,{}^B T_{TCP}
\end{equation}
接下来我们求 \({}^{TCP}T_O\)。我们已知\({}^W T_O\)，也就是\(O\rightarrow W\)。但目标不是世界坐标系，而是 TCP，因此还需要：\(W\rightarrow TCP\)。问题在于我们前面得到的是 \(TCP\rightarrow W\)。方向反了，必须取逆：\({}^{TCP}T_W=({}^W T_{TCP})^{-1}\)。因此完整路径 \(O\rightarrow W\rightarrow TCP\) 的表达式为
\begin{equation}
    {}^{TCP}T_O=({}^W T_{TCP})^{-1}\,{}^W T_O=\left({}^W T_B\,{}^B T_{TCP}\right)^{-1}{}^W T_O
\end{equation}


\subsection{特性参数}

特性参数（Kenngrößen）描述这台机器人到底能不能完成某项任务，分为4大类：
\begin{itemize}
    \item \textbf{几何参数}描述的是机器人的尺寸以及它能够覆盖的空间。
    \item \textbf{负载参数}描述的是机器人的负载。
    \item \textbf{精度参数}描述的是机器人操作的精度。包括定位精度和重复精度。
    \item \textbf{运动参数}描述的是例如循环时间（完成一套动作需要的时间）、速度、加速度、稳定时间等参数。
\end{itemize}
这四个参数共同描述了机器人的性能。

\subsection{末端执行器与安全技术}

决定机器人能干什么任务的是装在机械臂末端的\textbf{末端执行器（Endeffektor）}。

\paragraph{夹爪是怎么抓住零件的}

夹爪要把一个零件稳稳固定住，靠的无非是三种抓力原理。
\begin{itemize}
    \item \textbf{力锁合（Kraftschluss）}：靠某种持续作用的力把物体夹紧，比如靠摩擦力、负压或磁力。
    \item \textbf{形锁合（Formschluss）}：靠零件和夹具的形状互相咬合。
    \item \textbf{质锁合（Stoffschluss）}：靠物质本身的状态变化把物体粘住，比如冷冻。
\end{itemize}

这三种夹爪形态的示例如图 \ref{fig:halterprinzipien} 所示。力锁合最常见也最灵活，形锁合精度高但夹具要专门定制，质锁合则多用于特殊工艺（比如玻璃、易碎品的临时固定）。

\begin{figure}[htbp]
  \centering
  \includegraphics[width=0.8\textwidth]{halteprinzipien.png}
  \caption{三种夹爪形态的示例}
  \label{fig:halterprinzipien}
\end{figure}

光有夹爪还不够，实际生产中还需要以下这些装置来辅助完成任务：
\begin{itemize}
    \item 自动换爪系统：当机器人需要频繁切换不同的夹爪或工具时，可以在机器人手腕和工具之间加一块"适配板"，任务变了就自动换一套工具，不用人工去拧螺丝换夹爪。
    \item 碰撞保护装置：适配板之间靠气压支撑，正常工作时是刚性连接，一旦工具意外碰撞受到过载，气压结构会自动松开缓冲，避免损坏机器人或工件。
    \item 位置补偿装置：在装配这类需要高精度对准的工序中，补偿装置能允许末端执行器在小范围内自由浮动、自动纠正位置和角度偏差，同时又能防止绕装配轴转动跑偏，特别适合精密插装、螺丝拧入这类任务。
\end{itemize}

\paragraph{机器人周围的安全防护}

工厂里必须要有防护装置把人和机器人的工作区域隔开或及时停机。这些防护装置按方式可以分为两大类：
\begin{itemize}
    \item 隔离型防护：物理上把人挡在机器人工作区之外。比如围栏、安全门。
    \item 非隔离型防护：靠感知或操作方式来保证安全。比如安全垫（人一踩上去触发停机）、激光扫描仪（探测到人进入危险区并立即停机）、双手按钮开关（必须同时用两只手按下才能启动，防止手伸进危险区）。
\end{itemize}

\subsection{轨迹规划}

机器人要完成一个动作，在有了夹爪和骨架之后，还得进行\textbf{轨迹规划（Bahnplanung）}，它要做的就是从几何角度把一个运动任务描述清楚：路径约束、时间约束、机器人自身约束。

规划好的轨迹要变成机器人真正的运动，中间要经过几道工序，称为\textbf{面向指令段的轨迹控制（satzorientierten Bahnsteuerung）}。大致分四步：
\begin{enumerate}
    \item \textbf{为指令段做准备（Satzvorbereitung）}：把编程时的几何数据和运动学数据整理成数学函数
    \item \textbf{轨迹插补（Bahninterpolation）}：计算出轨迹上具体的点
    \item \textbf{坐标变换（Koordinatentransformation）}：把这些点从直角坐标系转换成机器人每个关节该转的角度值
    \item \textbf{位置控制（Lageregelung）}：把这些角度值作为目标值交给位置控制器去驱动电机实际执行。
\end{enumerate}

机器人的运动方式主要有三种。
\begin{itemize}
    \item \textbf{点位控制（PTP，Point to Point）}：只规定起点和终点，中间走的具体路径形状不确定，简单快速，常用于不要求路径形状的搬运任务。
    \item \textbf{多点控制（MP）}：通过预先采样、记录一系列密集的路径点来描述一条复杂曲线，常用示教/回放的方式编程。
    \item \textbf{连续轨迹控制（CP，Continuous Path）}：整条路径的位置和姿态在全程都被精确规定，常见的插补方式有直线、圆弧、样条曲线，其中直线轨迹最常用，比如装配这类需要精确路径的场合。这些插补方式如图 \ref{fig:interpolation} 所示。
\end{itemize}

\begin{figure}[htbp]
  \centering
  \includegraphics[width=0.35\textwidth]{interpolation.png}
  \caption{三种夹爪形态的示例}
  \label{fig:interpolation}
\end{figure}


\subsection{编程方法}

\subsubsection{在线编程}

在线编程的特点是直接在真实机器人上进行示教，主要有两种做法。
\begin{itemize}
    \item 示教编程（Teach-In）：如图 \ref{fig:teachin} 所示，操作人员用一个操作手柄，把机器人依次移动到任务需要的几个关键位置，每到一个位置就把该点的坐标连同相关参数（比如用哪种插补方式、走多快）一起存下来；正式运行程序时，机器人就按顺序依次开到这些存好的点，并套用当时记录的参数。
    \item 引导编程（Playback）：如图 \ref{fig:playback} 所示，不是一个点一个点地示教，而是直接拉着机器人的手走完整个动作，比如靠装在末端执行器上的力矩传感器感知人手的推力，或者靠一个运动学替代模型来跟随人的动作，机器人会把整个连续的运动过程记录下来，之后原样重放。
\end{itemize}
这两种方式的共同优点是直观、不需要专业编程知识，缺点是要占用真实机器人的时间，示教期间机器人没法干别的活。

\begin{figure}[htbp]
    \centering
    \begin{subfigure}[b]{0.46\textwidth}
        \centering
        \includegraphics[width=\textwidth]{teachin.png}
        \caption{示教编程 Teach-in}
        \label{fig:teachin}
    \end{subfigure}
    \hfill
    \begin{subfigure}[b]{0.49\textwidth}
        \centering
        \includegraphics[width=\textwidth]{playback.png}
        \caption{引导编程 Playback}
        \label{fig:playback}
    \end{subfigure}
    \caption{在线编程方法}
    \label{fig:onlineprogramming}
\end{figure}

\subsubsection{离线编程}

离线编程不需要占用产线上的真实机器人，如图 \ref{fig:offlineprogramming} 所示。

\begin{figure}[htbp]
  \centering
  \includegraphics[width=0.8\textwidth]{offlineprogramming.png}
  \caption{离线编程案例}
  \label{fig:offlineprogramming}
\end{figure}

离线编程也是有两种做法：
\begin{itemize}
    \item 文本编程：直接用文字代码写程序，因为只需要一个编辑器，投入成本很低；不过问题在于几乎每个机器人厂商都有自己的一套编程语言，彼此互不兼容。
    \item 基于仿真的编程：先在电脑里把整个加工单元（工作台、零件、机器人等）建模出来，规划好机器人大致怎么运作，再通过设置一系列途经姿态（Via-Pose）来编写机器人的运动路径。
\end{itemize}

仿真模型里规划的路径搬到真实产线上执行时，因为设备安装误差、零件实际位置和仿真设想的不完全一致等原因，机器人实际走的位置和姿态往往会和理论值有偏差（姿态误差）。为了解决这个问题，通常先在仿真环境里把整套系统规划出来，然后把真实设备搭建起来，接着对真实产线进行实测标定，最后根据标定结果反过来修正仿真模型，或者给机器人的原始轨迹加一个偏移量，让仿真规划出来的路径在真实世界里也能精确复现。


\subsubsection{混合编程}

在线和离线各有短板，混合编程就是把两者结合起来，先用示教在真实机器人上快速确定出几个关键位置，锚定住任务的大致框架，然后再用离线的方式去补充和精细化：一种做法是用示教编程定好关键点后，再用文本代码把中间的逻辑、参数补全；另一种做法是用示教编程定好关键点后，再放到仿真环境里做进一步的路径优化和碰撞检查。这样既保留了示教所见即所得、贴近真实情况的优点，又能用离线手段做更精细的编辑和验证，减少了纯离线编程常见的仿真与现实偏差问题。


\subsection{智能化融入}
\label{subsec:intelligent}

前面几节讲的都是机器人本身，它长什么样、怎么动、怎么编程。这一节换了个角度，讲的是怎么把机器人聪明地安排进一条有人参与的生产线，让人和机器人各自发挥所长、协同完成任务。

\paragraph{示例：飞机隔框铆接}

如图 \ref{fig:nietprozess} 所示，原来的做法是两名工人手持铆枪和顶把，把大量铆钉一颗颗打进隔框，把它和机身段连接起来。这个工作又累又不符合人体工学，因为它经常需要将手臂举过头顶作业，铆钉数量又多，是个很适合重新设计、引入自动化辅助的场景。

\begin{figure}[htbp]
  \centering
  \includegraphics[width=1\textwidth]{nietprozess.png}
  \caption{飞机隔框铆接的原始方法}
  \label{fig:nietprozess}
\end{figure}

要把人和机器人智能地组合到一条产线里，大致分六步：
\begin{enumerate}
    \item 流程规划与方案设计，明确任务该分给人还是机器人
    \item 建一个虚拟模型，把产品、工艺、设备信息都参数化地描述出来
    \item 做产线配置，把规划好的方案在现实中搭建起来并标定
    \item 设计，让工人能方便地了解和调整系统
    \item 实际操作运行
    \item 工艺执行
\end{enumerate}
这六步环环相扣，前一步的结果是后一步的输入。接下来将详细解释每一个步骤是如何实现的。

\paragraph{第一步：流程规划与方案设计}

方案设计的核心是能力匹配（Mapping）：先梳理出这项任务实际需要哪些能力，再对比人和机器人各自能提供哪些能力，把两者匹配起来，形成几种备选方案。总体来看，人的优势在于灵活，短板是对工艺的把控不够精确、容易受人体工学限制。机器人恰好相反：能严格按照设定路径精确移动、可靠地重复执行单调任务、结果一致性好；短板是产品设计必须迁就自动化要求，也不像人一样灵活应变。

放到铆接案例里，工人负责铆接任务里需要灵活操作、经验判断的部分，机器人则负责单调但需要稳定输出力的部分。所以外面的人操作铆枪，里面的机器人稳稳地顶住，两者配合完成铆接。

\paragraph{第二步：虚拟模型}

有了分工方案后，就可以搭建一个虚拟模型，参数化地记录。这是为了能提前在电脑里对方案做分析、验证和优化，减少后续返工的风险。

\paragraph{第三步：产线配置}

把实际需要的设备搭建到位，再通过测量确定它们真实的位置和朝向，这称之为\textbf{标定}，然后把这些实测数据反馈回虚拟模型里，让虚拟模型和现实情况对齐。机器人的实际控制就可以基于这个校准过的虚拟模型来进行。

标定本身还可以分三个层级，如图 \ref{fig:konfiguration} 所示：
\begin{figure}[htbp]
  \centering
  \includegraphics[width=0.8\textwidth]{konfiguration.png}
  \caption{标定的三个层级}
  \label{fig:konfiguration}
\end{figure}
\begin{itemize}
    \item \textbf{全局标定（Globales Einmessen）}：确定大件设备（比如工作台）在空间中的整体位置；
    \item \textbf{区域标定（Regionales Einmessen）}：修正某个工件本身摆放位置的偏差；
    \item \textbf{局部标定（Lokales Einmessen）}：进一步测量工件上具体特征点（比如某个铆钉孔）的精确偏差。
\end{itemize}
层级越细，精度要求也越高。案例中就用到了区域标定：用线激光三角测量传感器识别隔框上的铆钉孔图案，据此计算出一个局部坐标系，用来修正真实工件和CAD模型之间的偏差，进而推算出其余铆钉孔的准确位置，供机器人后续铆接使用。

\paragraph{第四步：人机交互界面}

产线上的工人往往不懂机器人编程，但生产中经常需要根据新的产品变型或流程优化去调整机器人程序。解决办法是给工人提供直观的操作界面和仿真环境，让他们不需要专业编程知识也能修改任务，调整结果还会实时反馈给工人确认。

\paragraph{第五步和第六步：实际操作运行工艺执行}

系统运行起来后，工人应当能在生产过程中随时进行简单的干预，比如用语音喊启动/停止来控制流程推进到下一步，不需要复杂培训就能上手，也让工人对整个流程更有掌控感、更容易接受这套人机协作方式。